%!TEX root = ../thesis.tex
% ===================================================================
% Chapter 3: System Design  (replaces the stubs in chapters/design.tex)
% Needs in imports.tex:
%   \usetikzlibrary{positioning,arrows.meta,shapes.geometric,fit,backgrounds,calc}
% Lines marked %CHECK need a value or label confirmed before you hand in.
% ===================================================================
\mychapter{System Design}{System Design}{}
\label{chap:design}

% -------------------------------------------------------------------
% Shared figure styles. Move into imports.tex if you prefer.
% -------------------------------------------------------------------
\tikzset{
  blk/.style={draw, rounded corners=2pt, align=center, font=\footnotesize,
              minimum width=2.6cm, minimum height=0.95cm, inner sep=3pt},
  mdl/.style={blk, fill=black!8},
  hw/.style={blk, dashed},
  st/.style={draw, rounded corners=2pt, align=center, font=\footnotesize,
             minimum width=2.7cm, minimum height=0.9cm},
  dec/.style={draw, diamond, aspect=2.2, align=center, font=\scriptsize,
              inner sep=1pt},
  arr/.style={-{Stealth[length=2mm]}, thick},
  lbl/.style={font=\scriptsize, fill=white, inner sep=1pt}
}

%% ===================================================================
\mysection{User and functional requirements}{User and Functional Requirements}
\label{sec:requirements}

The person I have in mind is someone in a home or small office who wants to
tell a rover what to do in ordinary speech, and who is not watching a screen
while it happens. They will say things like ``go and look at the orange box''
without checking whether there is one orange box or two. The rover has to cope
with that, and it has to be honest when it cannot. I turned this into the
requirements in Table~\ref{tab:requirements}. Each one traces to an objective in
Section~\ref{sec:objectives}, so the evaluation in Chapter~\ref{chap:evaluation}
can be read back against them.

\begin{table}[htbp]
  \centering
  \caption{Functional (FR) and non-functional (NR) requirements, with the objective each one serves.}
  \label{tab:requirements}
  \small
  \begin{tabular}{@{}l p{9.4cm} l@{}}
    \toprule
    ID & Requirement & Objective \\
    \midrule
    FR1 & Take the command by voice and transcribe it locally. & 2 \\
    FR2 & Before planning, check the command against the room and ask the person a spoken question when it is ambiguous, up to a fixed number of questions. & 4 \\
    FR3 & Produce a plan that uses only actions the rover has and objects the room contains. & 3, 5 \\
    FR4 & Read the plan back and move only after a spoken confirmation. A rejection must lead back to the dialogue, not end it. & 4 \\
    FR5 & Carry out the plan one step at a time, with obstacle stopping that does not depend on any model. & 1, 3 \\
    FR6 & Use the camera to check progress and arrival during a step. & 5 \\
    FR7 & Say at the end what happened and what the rover was still unsure about. & 4, 7 \\
    FR8 & When a step fails, ask the person what to do next and carry on from where the rover is. This is for failure only, never for ambiguity. & 4 \\
    FR9 & Log every model call, check and decision in a form the evaluation can count. & 6, 7, 8 \\
    FR10 & When a command contains a condition, look once to decide it, run only the matching steps, and ask the person when the camera cannot tell. Never guess a branch. & 3, 7 \\
    FR11 & When the person asks what is near a target, report what is around it and not only whether the target was reached. & 3, 5 \\
    NR1 & Run planning, vision and speech on a Ryzen 5 3600 with 16\,GB RAM and a 12\,GB GPU. Cloud models are allowed only as a benchmark. & 9 \\
    NR2 & Keep the rover stoppable when the network or a model is slow or gone. & 1 \\
    NR3 & Use only a camera, an ultrasonic sensor and a line sensor. No encoders, IMU or lidar. & 1 \\
    NR4 & Let the planner or vision model be swapped by configuration, with no change to the pipeline code. & 9 \\
    NR5 & Show the person what state the rover is in without a screen. & 4 \\
    \bottomrule
  \end{tabular}
\end{table}

NR3 is the one that shapes the rest of the chapter. Without odometry the rover
does not know how far it has moved. Every distance it covers is an open-loop
estimate, and the camera and sonar are what correct it.

The command levels in Table~\ref{tab:command-levels} each need something
different from the design, and Table~\ref{tab:level_map} says what. I made this
mapping before building the later parts, and it is the reason FR10 and FR11
exist. A level with no mechanism behind it would be a level I could not
honestly claim to test.

\begin{table}[htbp]
  \centering
  \caption{What the design provides for each command level.}
  \label{tab:level_map}
  \small
  \begin{tabular}{@{}l p{4.0cm} p{6.4cm}@{}}
    \toprule
    Level & Command type & Mechanism \\
    \midrule
    L1 & Single action & One \texttt{move} or \texttt{turn} step, read back and confirmed. \\
    L2 & Sequence of actions & Ordered steps, each finished before the next starts. \\
    L3 & Perception-dependent & \texttt{approach} to find and reach a named object (Section~\ref{ssec:freeroam}), or \texttt{follow\_line} to a visible end point. \\
    L4 & Conditional & A \texttt{check} step with tagged branches (Section~\ref{sec:conditional}). \\
    L5 & Vague & The clarification loop before departure (Section~\ref{sec:clarification_design}) and a survey of the surroundings when the command asks for one (Section~\ref{ssec:survey}). \\
    \bottomrule
  \end{tabular}
\end{table}

%% ===================================================================
\mysection{System architecture}{System Architecture}
\label{sec:architecture}

\mysubsection{Off-board inference and the thin-client split}{Off-board Inference}
\label{ssec:split}

The rover is a Hiwonder TurboPi with a Raspberry Pi 5. I run the speech
recogniser, the planner and the vision model on a separate machine, which I call
the brain, and keep only motion and sensing on the Pi. The reason is size. The
three models together take about 7\,GB, which does not fit on the Pi next to the
camera and drive stack, and a model small enough to fit would plan poorly. Running
them off-board also means I can swap a model by changing one environment
variable, which the cloud comparison in Chapter~\ref{chap:evaluation} depends on
(NR4).

The brain reaches the Pi over rosbridge on port 9090 using roslibpy
(Figure~\ref{fig:split}). Commands go one way. Camera frames, sonar readings and
step results come back the other. The brain owns the dialogue, the plan and the
mission loop. The Pi owns whatever has to keep working when the link is slow,
which is the sonar stop and the line follower. If the network drops during a
step, the rover still stops for an obstacle (NR2).

\begin{figure}[htbp]
  \centering
  \begin{tikzpicture}[scale=0.85, transform shape]
    \node[mdl] (stt) at (0,1.2)  {Speech\\recogniser};
    \node[mdl] (pln) at (3.0,1.2){Planner\\(LLM)};
    \node[mdl] (vlm) at (6.0,1.2){Vision\\model};
    \node[blk] (dlg) at (1.5,-0.3){Dialogue and\\mission loop};
    \node[blk] (tts) at (4.5,-0.3){Speech\\output};
    \begin{scope}[on background layer]
      \node[draw, rounded corners=3pt, fit=(stt)(pln)(vlm)(dlg)(tts),
            inner sep=7pt, label={[font=\footnotesize]above:Brain (desktop or laptop)}] (brain) {};
    \end{scope}
    \node[hw] (react) at (11.2,1.2) {Reactive layer\\sonar stop, line\\follower};
    \node[hw] (io)    at (11.2,-0.3){Camera, gimbal,\\wheels, LEDs};
    \begin{scope}[on background layer]
      \node[draw, rounded corners=3pt, fit=(react)(io), inner sep=7pt,
            label={[font=\footnotesize]above:Rover (Raspberry Pi 5)}] (pi) {};
    \end{scope}
    \draw[arr, <->] (brain.east) -- node[lbl, above]{rosbridge} (pi.west);
  \end{tikzpicture}
  \caption{The compute split. The three models and the dialogue run on the brain. The Pi runs only what must keep working without them.}
  \label{fig:split}
\end{figure}

The code boundary is one method, \texttt{RoverController.execute\_step}. The brain
hands it one step and gets back \texttt{ok}, \texttt{blocked} or \texttt{halted}.
Everything above that method is the same for the real rover and for a simulated
one, so I can run the same dialogue and mission code in simulation before
putting it on hardware.

The default configuration is fully local. Speech recognition is
\texttt{faster-whisper} with the \texttt{fluister-turbo} model, the planner is
\texttt{gemma4:e2b}, and the vision model is \texttt{qwen2.5vl:3b}. Both are
served by Ollama with a context of 8192 tokens. Speech output uses Piper with the
\texttt{en\_US-lessac-medium} voice. Cloud providers go through the same
interface, and I use them only as a comparison in
Section~\ref{sec:cloud_comparison}.

The 8192-token window is a constraint I check instead of assuming. Every
clarification turn and every recovery turn (Section~\ref{sec:recovery}) adds text
to a prompt, and a small model that is silently truncated plans worse without
raising any error. A separate offline tool, \texttt{context\_budget.py}, counts
the tokens each part of the prompt takes (fixed instructions, room digest, one
clarification turn, one recovery turn) so I know how many turns fit before the
window is the limiting factor. The measured counts are in
Chapter~\ref{chap:evaluation}.

\mysubsection{Two-rate control: waypoint intents versus the reactive layer}{Two-rate Control}
\label{ssec:rates}

A vision-model call takes about 15\,s on my hardware. Keeping a mecanum chassis
on a line needs a command every 30 to 100\,ms. A steering loop cannot go through
a model that slow, so I split control into two layers that run at different rates.

The reactive layer runs on the Pi at 10 to 30\,Hz. It reads the sonar, zeroes the
wheels when something is in the way, and follows the line from the IR array
without asking a model anything. The deliberative layer runs on the brain at the
pace of the models. It chooses the next step, decides whether the rover has
arrived, and looks at a frame when a step finishes. The two layers meet at step
boundaries and at a small number of checks during a step, not at every control
tick.

The consequence I care about is that the vision model never sits in a safety
path. A slow or wrong answer from it can send the rover to the wrong place, which
I measure, but it cannot make the rover drive into a wall, because the stop is on
the Pi.

\mysubsection{The end-to-end pipeline}{Pipeline}
\label{ssec:pipeline}

Figure~\ref{fig:pipeline} shows the path of one command. Before any command, I
record a short walkthrough video of the room. The brain takes up to ten
keyframes from it, has the vision model list the objects in each, and condenses
the result into a short text description that I call the scene digest. The digest
is built once and reused. It is capped at 8 objects with 2 attributes each, so it
fits the planner prompt without crowding out the dialogue. I chose text over a
map deliberately. A map is a separate system to build and keep right, and
building one is outside the scope set in Section~\ref{sec:scope-and-limitations}.

A spoken command is transcribed and enters the clarification loop
(Section~\ref{sec:clarification_design}). When the command is clear, the planner
turns it into steps using the vocabulary of the active rover profile
(Section~\ref{sec:vocabulary}). The vision model then checks the plan against the
digest for steps that cannot work, for example a target that is not in the room.
The plan is read back, and the rover moves only after a spoken confirmation. The
mission loop then runs the steps one at a time, and at the end a report is logged
and spoken.

Two WebSocket endpoints carry this. \texttt{/ws/dialogue} serves everything up to
confirmation, and \texttt{/ws/execution} serves the mission. I kept them apart so
the dialogue state and the mission state never share a session object. A dialogue
session that sits idle for 30 minutes is dropped.

\begin{figure}[htbp]
  \centering
  \begin{tikzpicture}[scale=0.82, transform shape]
    \node[mdl] (scene) at (7.0,2.4) {Scene digest\\(video $\rightarrow$ keyframes\\$\rightarrow$ vision model $\rightarrow$ text)};
    \node[blk] (voice) at (0,0)    {Spoken\\command};
    \node[mdl] (stt)   at (3.5,0)  {Speech\\recognition};
    \node[mdl] (clar)  at (7.0,0)  {Ambiguity check\\(asks until clear)};
    \node[mdl] (plan)  at (10.5,0) {Planner};
    \node[mdl] (ver)   at (10.5,-2){Plan\\check};
    \node[blk] (conf)  at (7.0,-2) {Read-back and\\confirmation};
    \node[hw]  (exec)  at (3.5,-2) {Mission loop\\(rover steps)};
    \node[blk] (rep)   at (0,-2)   {Report\\(spoken, logged)};
    \draw[arr] (voice)--(stt);
    \draw[arr] (stt)--(clar);
    \draw[arr] (clar)--node[lbl]{clear}(plan);
    \draw[arr] (plan)--(ver);
    \draw[arr] (ver)--(conf);
    \draw[arr] (conf)--node[lbl]{yes}(exec);
    \draw[arr] (exec)--(rep);
    \draw[arr, dashed] (scene.south) -- (clar.north);
    \draw[arr, dashed] (scene.east) -| (12.3,-2) -- (ver.east);
  \end{tikzpicture}
  \caption{One command from speech to report. Shaded blocks call a model. The dashed block drives the rover. Dashed arrows show the scene digest feeding the two checks that need to know what is in the room.}
  \label{fig:pipeline}
\end{figure}

\mysubsection{Status lights}{Status Lights}
\label{ssec:lights}

NR5 asks that the person can tell what the rover is doing without looking at a
screen. The Pi has RGB LEDs on the board and on the sonar ring, so I use them as
a state display (Table~\ref{tab:lights}). The brain sets a light at each phase
change through one call, \texttt{status\_light}, which returns at once and never
raises. A failed light must not be able to stop a mission. Rovers that have no
LEDs inherit a no-op, and one setting (\texttt{RGB\_ENABLED}) turns the whole
thing off.

\begin{table}[htbp]
  \centering
  \caption{The eight light states.}
  \label{tab:lights}
  \small
  \begin{tabular}{@{}l l p{6.6cm}@{}}
    \toprule
    State & Pattern & Meaning \\
    \midrule
    listening & steady blue & Waiting for a command or an answer. \\
    thinking & red and green alternating, 0.5\,s & A model is working. \\
    confirm & steady amber & A plan has been read back and the rover is waiting for yes or no. \\
    executing & steady green & Carrying out a step. \\
    searching & steady cyan & Looking for a target. \\
    awaiting\_guidance & magenta, 1\,s on and off & Stopped and waiting for the person to say what to do next. \\
    arrived & three green flashes, then steady green & A target has been reached. \\
    blocked & steady red & The mission has stopped and cannot continue. \\
    \bottomrule
  \end{tabular}
\end{table}

The mission loop logs each change with a timestamp, so the same signal that the
person sees is also in the run log and can be drawn as a timeline afterwards.

%% ===================================================================
\mysection{The capability vocabulary}{The Capability Vocabulary}
\label{sec:vocabulary}

A planner that may write any step it likes will write steps the rover cannot do.
A command like ``bring me the bottle'' invites a \texttt{pick\_up} step, and the
TurboPi has no arm. I give the planner a fixed list of actions and say in the
prompt which exist and what each one means. The
list depends on the rover, so it lives in a planner profile, one file per rover,
chosen by the \texttt{ROVER} variable or overridden with
\texttt{PLANNER\_PROFILE}.

On the real rover the base actions are \texttt{follow\_line}, \texttt{move},
\texttt{turn}, \texttt{approach}, \texttt{scan}, \texttt{observe}, \texttt{stop}
and \texttt{report}. A step is a small JSON object with an id, an action and a
target. The simulated rover's profile is the same except that
\texttt{follow\_line} is absent and \texttt{move} and \texttt{turn} carry a
distance and an angle, because the simulator can execute an exact value and the
Pi cannot (NR3). The prompt also states that the rover has no arm and cannot pick
up or carry anything, so a request for it is answered with what the rover can do
instead.

A ninth action, \texttt{check}, is not always on the list. It means ``look once
and decide whether a stated condition is true'' (Section~\ref{sec:conditional}),
and I offer it only when the command contains one of the words \emph{if},
\emph{whether}, \emph{otherwise} or \emph{unless}. In that case the prompt also
gains a short block that explains the \texttt{check} step and the optional
\texttt{when} field a step then carries (\texttt{yes} or \texttt{no}), with one
worked example. I left \texttt{check} out of every other prompt for two reasons.
Every extra line is context the 8192-token window has to pay for, and an action
the planner has no use for is one it can misuse, for instance by inventing a
condition nobody spoke. The cost of deciding this with a word match is that a
conditional phrased in other words, such as ``when you see the door, stop'', gets
no \texttt{check} action and the planner cannot express it. I accept that limit
for now and name it in the evaluation. Conditionals exist only on the real rover's
profile, not the simulated one.

I should be exact about where this constraint is enforced, because it is easy to
overstate. For the first plan it lives in the prompt. Nothing between the planner
and the read-back rejects a step whose action is outside the list. Three things
catch it afterwards. A step with an unknown action returns \texttt{blocked} when
the rover tries to run it. The end-of-mission report counts the out-of-vocabulary
actions in each plan, so I can measure how often the planner ignores the list
(Section~\ref{sec:executability}). And the person hears the plan before anything
moves. I chose to measure the first-plan failure instead of silently repairing it,
because how often the planner leaves the vocabulary is itself a result for RQ1.
The one repair I do make is on the real rover in free-roam mode, where a
\texttt{move} step that names an object is rewritten to \texttt{approach}, since
the planner often reaches for the wrong word and the intent is clear.

A plan produced during recovery is different, and there I do enforce the list.
When a step fails and the person gives new instructions, the replacement plan
goes straight into the mission with no spoken read-back, so nobody has heard it.
\texttt{validate\_steps} therefore checks that every step is an object whose
action is in the profile, and a plan that fails the check is thrown away
(Section~\ref{sec:recovery}). In free-roam mode the list for a replan leaves out
\texttt{follow\_line}, since there is no floor line for it to follow.

Keeping the profile in its own file has one more benefit. A hash of the profile's
vocabulary and step format is logged with every run, so when I change the wording
I can tell which results came from which version. The hash covers the base
profile and not the per-command variant, so it does not separate a conditional
prompt from a normal one. The command is in the log, which tells me which variant
was used. %CHECK: say this only if you leave profile_info() as it is.

%% ===================================================================
\mysection{Clarification loop design}{Clarification Loop Design}
\label{sec:clarification_design}

\mysubsection{Where the loop sits}{Placement}

Spoken commands are often underspecified, and the cost of acting on a wrong guess
is a rover in the wrong place. I handle ambiguity at three moments, and this
section covers the first. Before departure, the system asks when the command is
unclear. During the mission it does not stop to ask about ambiguity. It compares
what it sees with the digest and corrects itself if it can
(Section~\ref{ssec:revision}). At the end, it says what it was still unsure
about. I leave out a mid-mission question about ambiguity on purpose. The person
is often not next to the rover, and a robot that stops and waits for an answer it
may not get is worse than one that makes its best attempt and says what it did.
This differs from KnowNo~\cite{ren2023knowno} and Introspective
Planning~\cite{liang2024introspective}, which ask for help when their confidence
drops, and I come back to the difference in Section~\ref{sec:clarification_lit}.

\mysubsection{The dialogue as a state machine}{State Machine}
\label{ssec:statemachine}

A command moves through the phases in Figure~\ref{fig:states}. In
\texttt{CLARIFYING}, the vision model is asked whether the command is ambiguous,
given the command, the dialogue so far and the scene digest. It answers in JSON
with four fields: whether it is ambiguous, why, a question to ask, and a rewritten
command. If it is ambiguous, the question is spoken and the answer is added to the
history. If it is clear, the rewritten command is used from then on. I require
that this command names the object only and avoids words like ``left'' or
``behind you'', because those are relative to the camera view at the time of
asking and are wrong once the rover has moved.

The loop is capped at five questions. If the command is still unclear after that,
the system marks the session as capped and plans from the rewritten command if
the model produced one, and from the original command if it did not. In that case
it no longer appends the dialogue so far to the command. %CHECK: add your reason for this change.
The loop also drops a question it has already asked. If the only question the
model offers is one that is already in the dialogue, there is nothing new to ask,
so the system treats the session the same way as a capped one and moves on to
planning with the warning in the read-back. The read-back
then begins with ``I am not fully sure I understood'', so the person knows the
plan rests on a guess. I chose a cap over an open loop because I would rather have
a person confirm or reject a flawed plan than get tired of questions and walk
away.

A failed vision call is never read as ``nothing is unclear''. If the call errors,
I retry once. If it fails again, the person is told and asked to repeat the
command. Without this, an unavailable model would look the same as a clear
command and let everything through.

\begin{figure}[htbp]
  \centering
  \begin{tikzpicture}[scale=0.9, transform shape]
    \node[st] (cl) at (0,0)    {CLARIFYING};
    \node[st] (pl) at (4.2,0)  {PLANNING};
    \node[st] (ve) at (8.4,0)  {VERIFYING};
    \node[st] (aw) at (4.2,-2.6){AWAITING\_\\CONFIRMATION};
    \node[st] (ca) at (8.4,-2.6){CANCELLED};
    \node[st] (ex) at (4.2,-5.0){EXECUTING};
    \node[st] (re) at (8.4,-5.0){REPORTING};
    \draw[arr] (cl)--node[lbl]{clear}(pl);
    \path (cl) edge[arr, loop above, min distance=7mm]
          node[lbl, above]{ambiguous: ask} (cl);
    \draw[arr] (pl)--(ve);
    \draw[arr] (ve.north) to[out=120,in=60]
          node[lbl, above]{unworkable: replan once} (pl.north);
    \draw[arr] (ve.south west)--node[lbl, right]{plan ready}(aw.north east);
    \draw[arr] (aw)--node[lbl, right]{yes}(ex);
    \draw[arr] (ex)--(re);
    \draw[arr] (aw.west) -| node[lbl, pos=0.3, below]{no: ``What should I change?''} (cl.south);
    \draw[arr] (aw)--node[lbl]{cancel}(ca);
  \end{tikzpicture}
  \caption{The dialogue phases up to the start of a mission. A rejection returns to \texttt{CLARIFYING} with the person's reason added to the history.}
  \label{fig:states}
\end{figure}
% CHECK: CANCELLED is drawn only from AWAITING_CONFIRMATION. Redraw if other phases can cancel.

\mysubsection{Checking the plan before reading it back}{Plan Check}
\label{ssec:plancheck}

After the planner produces a plan, the vision model checks it against the digest
and returns whether it is workable and a list of concerns. If it says no, the
planner runs once more with the concerns added as advice while the original goal
stays as it was. I replan only once. A second round rarely fixes what the first
did not, and each round costs another model call. If the checker itself is
unavailable, the answer is neither yes nor no, and I do not replan. Skipping the
check is a safer failure than inventing a concern.

\mysubsection{Confirmation}{Confirmation}
\label{ssec:confirmation}

The plan is read back as a short sentence per step, and the rover does not move
until the person answers. The answer is classified in two stages. A keyword pass
runs first over English and Afrikaans phrase lists, with the reject list checked
before the confirm list, so ``no, don't go'' is not taken as a yes because it
contains ``go''. If the keywords find nothing, the planner model classifies the
reply. If that also fails or is unsure, the answer counts as unclear and the rover
stays where it is. I made the safe outcome the default, since a missed ``yes''
costs a few seconds and a missed ``no'' costs a moving robot. Unclear answers are
retried twice.

A rejection does not end the session. The system asks ``What should I change?''
and returns to \texttt{CLARIFYING} with the answer in the history, so the person
can fix one step without starting over.

Commands are accepted by voice only. A typed route exists for testing, but it is
off unless \texttt{ALLOW\_TEXT\_COMMANDS} is set, so the dialogue I evaluate is
the one a person would actually use.

\mysubsection{The two later moments}{Later Moments}
\label{ssec:revision}

Once the rover is moving, the second moment is self-correction. While a step
runs, new frames are compared with the digest. A small module, \texttt{revision},
classifies what it sees as no change, a reroute, a material change or a block,
and it does this with plain code and no model call. The planner may propose a
change, but the classifier decides whether the person has to be asked. A reroute
is applied between steps. A material change moves the mission into
\texttt{AWAITING\_REVISION\_CONFIRMATION}, and the rover waits for a spoken yes
before it follows the new plan. The third moment is the report
(Section~\ref{sec:report}).

%% ===================================================================
\mysection{Failure recovery}{Failure Recovery}
\label{sec:recovery}

Some things go wrong that no amount of questioning beforehand would have
prevented. The orange box is behind a chair. Someone has left a bag in the
corridor. The rover reaches the place the plan named and the camera does not
agree. In these cases the plan was clear and the world did not cooperate.

When a step returns \texttt{blocked}, the mission enters
\texttt{AWAITING\_GUIDANCE} (Figure~\ref{fig:guidance}). The rover says why it
stopped, using one of four fixed sentences chosen by the reason (target not found,
obstacle, arrival contradicted, or out of time) or a general sentence for
anything else, and asks what to do next. The
light turns magenta. The person answers by voice. The planner then writes a new
plan for what is still left, given the original command, the steps already done,
the step that failed and why, what the rover has recently seen, and the
conversation so far. The prompt tells the planner to use the answer to decide
where to look or what to reach instead. If the person says where the target is,
the plan should turn that way and approach the target again, and it must still end
by reaching the target. The output then goes through two checks. The first is
\texttt{validate\_steps} (Section~\ref{sec:vocabulary}). The second applies when
the failed step was an \texttt{approach}: a plan that never approaches anything is
rejected, because a plan that walks away from the failed step is not a recovery.
If the plan passes both, the mission carries on from where the rover is, and the
rover says ``Okay. Trying again.''

\begin{figure}[htbp]
  \centering
  \begin{tikzpicture}
    \node[st] (bl) at (0,0)    {Step returns\\\texttt{blocked}};
    \node[st] (aw) at (4.0,0)  {Speak reason,\\ask what next};
    \node[dec] (an) at (8.0,0) {answer?};
    \node[st] (rp) at (8.0,-2.2){Replan from\\current state};
    \node[dec] (ok) at (4.0,-2.2){steps pass\\\texttt{validate\_steps}?};
    \node[st] (go) at (0,-2.2) {Resume mission\\(EXECUTING)};
    \node[st] (hl) at (8.0,2.2) {HALTED\\``I cannot continue''};
    \draw[arr] (bl)--(aw);
    \draw[arr] (aw)--(an);
    \draw[arr] (an)--node[lbl, right]{an instruction}(rp);
    \draw[arr] (rp)--(ok);
    \draw[arr] (ok)--node[lbl]{yes}(go);
    \draw[arr] (ok.north) -- node[lbl, left]{no: ask again}(aw.south);
    \draw[arr] (an)--node[lbl, right]{stop, timeout, or 3 turns used}(hl);
  \end{tikzpicture}
  \caption{The recovery loop. The turn counter is shared across the whole mission, so it cannot be used to drag a failing mission on.}
  \label{fig:guidance}
\end{figure}

Three limits keep this from becoming an open-ended chat. There are at most three
guidance turns per mission. A reply has 120\,s to arrive. And the person can
decline. A reply that begins with ``stop'', ``cancel'', ``give up'' or ``never
mind'', or a bare short ``no'', ends the mission with the rover halted. I did not
reuse the confirmation classifier for this. Its reject list contains ``no'',
``wait'' and ``wrong'', so ``there is no box on the left'' would end the mission,
and that is exactly the kind of answer I am asking for.

Two decisions here are easy to miss. First, recovery fires only for a physical or
search failure. It never fires for ambiguity. Ambiguity is asked about before
departure and reported at the end, and keeping recovery separate from it is what
stops the design from collapsing into KnowNo. Second, a run that needed an answer
is never logged as a plain success. The report classes a mission as completed,
completed after guidance, halted after guidance, halted or aborted, so the
evaluation can separate what the rover did alone from what it did with help.
The report also carries a separate \texttt{recovered} flag, and it is stricter. It
is set only if an \texttt{approach} step succeeded after the last blocked step. A
mission that was talked through and then simply ran out of plan is not a
recovery, and I do not want it counted as one.

I do not call \texttt{halt()} on the way into recovery. A blocked step has already
ended its own timed motion, and on the Pi \texttt{halt()} sets a flag that only
the next \texttt{execute\_step} clears. The free-roam approach in
Section~\ref{ssec:freeroam} calls the drive primitives directly, so a latched flag
would make every later hop fail.

%% ===================================================================
\mysection{Conditional commands}{Conditional Commands}
\label{sec:conditional}

A conditional command, level L4 in Table~\ref{tab:command-levels}, asks the rover
to decide something and act on the answer: ``check whether the orange box is
visible, and if it is, go to it, otherwise stop''. The design question is who
makes the decision and when. I split it in two. The planner writes both branches
up front, before anything moves, and the rover only evaluates the one condition at
the moment it matters. That keeps the planner doing the thing it is good at,
turning words into a structure, and leaves the single judgement that needs a
camera to the vision model.

Table~\ref{tab:cond_plan} shows the plan shape the prompt asks for. One
\texttt{check} step states the condition as a short sentence. Steps for the yes
case carry \texttt{when: yes}, steps for the no case carry \texttt{when: no}, and
steps with no tag always run.

\begin{table}[htbp]
  \centering
  \caption{A conditional plan, as the prompt shows it to the planner.}
  \label{tab:cond_plan}
  \small
  \begin{tabular}{@{}l l l l@{}}
    \toprule
    Id & Action & Target & When \\
    \midrule
    1 & check & the orange box is visible & \\
    2 & approach & orange box & yes \\
    3 & report & orange box & yes \\
    4 & turn & left & no \\
    5 & stop & & no \\
    \bottomrule
  \end{tabular}
\end{table}

When the mission loop reaches the \texttt{check} step, the rover centres the
gimbal, waits half a second for it to settle, and takes one frame
(Figure~\ref{fig:cond}). The vision model is asked whether the stated condition is
true of what is visible and may answer yes, no or unsure. The loop remembers the
answer. Every later step tagged with the other answer is skipped and logged as
skipped, so the run record shows which branch ran and which did not. Replacing the
plan, which recovery can do, clears the remembered answer.

\begin{figure}[htbp]
  \centering
  \begin{tikzpicture}[scale=0.9, transform shape]
    \node[st] (ck) at (0,0)     {\texttt{check} step:\\look once, centred};
    \node[dec] (an) at (4.4,0)  {answer?};
    \node[st] (ys) at (8.6,1.5) {Run steps tagged\\\texttt{yes}, skip \texttt{no}};
    \node[st] (no) at (8.6,0)   {Run steps tagged\\\texttt{no}, skip \texttt{yes}};
    \node[st] (bl) at (8.6,-1.7){\texttt{blocked}:\\ask the person};
    \draw[arr] (ck)--(an);
    \draw[arr] (an)--node[lbl, above]{yes}(ys.west);
    \draw[arr] (an)--node[lbl, above]{no}(no.west);
    \draw[arr] (an)--node[lbl, below]{unsure or error}(bl.west);
  \end{tikzpicture}
  \caption{How a condition is resolved. An answer the camera cannot support does not pick a branch. It stops the step and goes to the recovery loop.}
  \label{fig:cond}
\end{figure}

The important choice is what happens on ``unsure'', and on a failed vision call,
which I treat the same way. The step returns \texttt{blocked}, and the existing
recovery loop (Section~\ref{sec:recovery}) asks the person what to do. I did not
add a default branch. If the rover guessed ``no'' whenever it could not tell, a
wrong guess would send it the wrong way and nobody would know it had been a guess.
Reusing the recovery loop also means conditionals need no second mechanism for
asking for help.

There are limits, and I would rather state them than have them found. The check is
one frame from wherever the rover is standing. It does not search. A condition
about an object just outside the camera's view will come back ``no'' or
``unsure'' even if the object is in the room, so the result depends on where the
rover was facing. The remembered answer is one value, so conditions cannot be
nested, and a tagged step that runs before any \texttt{check} has run does so
unconditionally. The spoken read-back lists each step as an action and a target
without its tag, so the person hears every branch as if it would run. %CHECK: fix the read-back (summarise_plan) or leave this sentence in.

%% ===================================================================
\mysection{Navigation as instrumentation}{Navigation as Instrumentation}
\label{sec:navigation}

Navigation is not what this project is about, but the pipeline needs the rover to
go somewhere so that a plan can succeed or fail on something real. I built two
ways of getting there and kept both simple. One follows a taped line. The other
finds a named object with the camera and drives to it.

\mysubsection{Free-roam approach to a named object}{Free-roam Approach}
\label{ssec:freeroam}

For \texttt{approach} the rover has no map and no odometry (NR3), so the vision
model finds the target, and the sonar and a second vision call decide whether it
has arrived. Figure~\ref{fig:approach} shows the loop.

\begin{figure}[htbp]
  \centering
  \begin{tikzpicture}[scale=0.9, transform shape]
    \node[st] (s) at (0,0)      {Search:\\pan the gimbal};
    \node[st] (l) at (4.0,0)    {Lock: one pivot,\\then verify};
    \node[st] (h) at (8.0,0)    {Hop toward\\target};
    \node[st] (b) at (0,-2.6)   {\texttt{blocked}\\with a reason};
    \node[dec] (a) at (4.0,-2.6){arrived?};
    \node[st] (c) at (8.0,-2.6) {Look again,\\gate the box};
    \node[st] (o) at (4.0,-4.8) {Second look\\cross-check};
    \node[st] (d) at (8.0,-4.8) {\texttt{ok}};
    \draw[arr] (s)--(l);
    \draw[arr] (l)--(h);
    \draw[arr] (h)--(c);
    \draw[arr] (c)--(a);
    \draw[arr] (a)--node[lbl, right]{yes}(o);
    \draw[arr] (o)--node[lbl]{confirmed}(d);
    \draw[arr] (a.north east)--node[lbl, above left]{no}(h.south west);
    \draw[arr] (s)--node[lbl, right]{not found}(b);
    \draw[arr] (a)--node[lbl, below]{obstacle, budget}(b);
    \draw[arr] (o.west) -| node[lbl, pos=0.2, above]{contradicted} (b.south);
  \end{tikzpicture}
  \caption{The free-roam approach. Each cycle is one hop and one look. Every exit that is not \texttt{ok} returns \texttt{blocked} with a reason, and that is what starts recovery.}
  \label{fig:approach}
\end{figure}

Search is done with the gimbal, not by turning the chassis. The camera pans
through a fixed list of angles, at most 18 vision calls in three sweeps, and the
chassis pivots by 110\textdegree{} between sweeps only if the first sweep finds
nothing. I chose gimbal-first because a pan is quick and its angle is known from
the servo setting, while a pivot on a mecanum base is an open-loop turn that
drifts. When the vision model reports the target at some horizontal position in
the frame, the bearing is the pan angle plus the offset from the frame centre
scaled by the measured 56\textdegree{} field of view. The rover then makes one
pivot toward that bearing and looks again to verify.

The vision model is a 3-billion-parameter model and it makes mistakes about
boxes. I do not take its bounding box at face value. Every detection passes a
plausibility gate that rejects a box that fills more than 60\% of the frame, a box
over 95\% of the frame's width or height, a box too small to be real, a box whose
area jumps threefold between hops while the rover is still far away, and a large
box sitting exactly at the centre. A detection that the model labels low
confidence is held until a second look at the same pan angle puts the box within
0.15 of the frame width of the first. Frames below a sharpness threshold are
never sent to the model, since a blurred frame gives it nothing reliable to
work from. The rover retries the grab twice and then skips the look.

Distance is where odometry would normally help. Lacking it, I scale each hop by
what the camera and sonar say. The default is a fixed pair: a 50\,cm hop while the
target looks small in the frame, and a 20\,cm hop once it fills more than 8\% of
it. A variable mode asks the vision model for a distance estimate in the same
call that locates the target, so it costs no extra call, and hops 60\% of the
estimated distance minus a 30\,cm stop margin, clamped between 8 and 50\,cm. If
the target is centred, the sonar clearance can shorten that further. I use a
fraction of the estimate, not all of it, because the number is uncalibrated and
open-loop hops drift. If the model returns no usable distance, the hop falls back
to the fixed pair. I built both modes because I do not know yet which is better,
and Chapter~\ref{chap:evaluation} compares them. The fixed mode stays the default
until that result is in.

Arrival is decided by two independent signals. The sonar says arrived when it
reads 300\,mm or less with the target centred. The camera says arrived when the
box fills 20\% of the frame or its bottom edge reaches 90\% of the frame height,
but only after at least one hop, only if the box did not shrink, and only if the
sonar does not read more than 800\,mm. Either signal then triggers a second vision
call that asks whether the target is still visible. If the answer is no, the rover
steps back, looks again, and returns \texttt{blocked} with the reason
\texttt{arrival\_contradicted} if the target stays hidden. I added this check
because a single positive from a 3B model, or from a sonar reading against some
other object, is not enough to call a step done.

The sonar threshold deserves a note. A sonar stop in the middle of a hop is common
at 300\,mm, because the hop ends inside the stopping distance. If the last look
showed the target centred, the first such stop is treated as the end of the hop and
the loop looks again. A second consecutive sonar stop without arrival is treated as
an obstacle and the step returns \texttt{blocked}. %CHECK ROVER_PI_SONAR_STOP_MM in .env was still 200 at last read; the code default is 300.

The loop has four ways to end in \texttt{blocked}: the target is not found after
the sweeps, an obstacle stops the second hop, the arrival cross-check
contradicts, or the cycle budget runs out (20 cycles or 240\,s). Each is the
reason the rover speaks in Section~\ref{sec:recovery}.

\mysubsection{Looking around: scan and survey}{Scan and Survey}
\label{ssec:survey}

A \texttt{scan} step on its own only moves the gimbal, which sees nothing. So
that ``tell me what you see'' has an answer in the log, the mission loop stops the
gimbal at three positions, left, right and centre, takes a frame at each, and has
the vision model describe it in one sentence. The rover is stationary the whole
time. That respects the control split in Section~\ref{ssec:rates}, since the
vision model never runs while the rover is moving. The gimbal finishes facing
forward.

Level L5 asks for more than a sentence. A command like ``what is near the orange
box'' is about an area, and a one-line description loses the small things. When
the command contains a word such as \emph{near}, \emph{nearby}, \emph{around},
\emph{beside}, \emph{next to} or \emph{everything}, the loop runs the same three
looks after the \texttt{observe} or \texttt{report} step, but asks a different
question: list every object in the frame, with its colour and where it sits,
each one once and at most ten. I limited the list because a small model asked to
list everything tends to repeat itself, which I had already seen when building
the room digest. The survey runs at most once per mission.

What the rover saw goes into the report as a list per direction. The spoken
summary is then told to name every listed object, grouped by left, centre and
right, and is allowed up to six sentences for that case instead of three. The
limits are plain. The rover looks in three directions and never behind itself,
and an object the model does not list is an object the report cannot mention.

\mysubsection{The line follower as the reactive layer}{The Line Follower}
\label{ssec:line_follower}
% KEEP your existing text for this subsection here, unchanged, and add this paragraph at its end:
%
% If the node's services cannot be reached, the supervisor now returns
% \texttt{blocked} with the reason \texttt{line\_follow\_unavailable} instead of
% letting the rover appear to run. Without that check, a missing node looked like
% a rover that had been told to follow a line and had simply not moved.
%CHECK: edit the last sentence if you do not want to describe the failure that way.

%% ===================================================================
\mysection{Reporting what happened}{Reporting}
\label{sec:report}

The third ambiguity moment is the report at the end of every mission, whether it
completed, was halted or was aborted. It has two parts. A block of facts is built
by plain code from the run log: which steps ran, how long each took, how many
model calls there were, what the approach did, and which models were used. Rules
then flag residual uncertainty from the same log. Examples are a target that was
not visible at arrival, a vision arrival that the sonar contradicted, an approach
that overran its time budget, and an action outside the vocabulary. Only the
spoken summary uses a model call, and it is given the facts and the flags so it
has nothing to invent. If that call fails, a fixed summary is spoken instead and
the facts are still logged.

The same log feeds Chapter~\ref{chap:evaluation}, which is why FR9 is a
requirement and not an afterthought. A mission is written to disk as a single
record that includes the plan, the checks, the guidance turns, the light timeline
and the model tag, so a run can be replayed and counted without the dashboard. Each condition that
was checked is listed with its answer, and when the person asked about
surroundings the record lists what was seen in each direction. A test level can
be tagged on a batch of runs, so the evaluation can group them without reading
every command.

%% ===================================================================
\mysection{Summary of design decisions}{Summary of Design Decisions}
\label{sec:design_decisions}

Table~\ref{tab:decisions} lists the decisions that shape the rest of the report,
with the reason for each and the alternative I turned down.

\begin{table}[htbp]
  \centering
  \caption{The main design decisions.}
  \label{tab:decisions}
  \small
  \begin{tabular}{@{}p{3.6cm} p{5.6cm} p{4.0cm}@{}}
    \toprule
    Decision & Reason & Turned down \\
    \midrule
    Models on a separate machine & About 7\,GB of models do not fit on the Pi. & Everything on the Pi, with a much smaller model. \\
    Reactive layer stays on the Pi & A model call takes 15\,s, steering needs 30 to 100\,ms. & Closing the loop through the vision model. \\
    Room as a text digest & A map is a separate system and out of scope. & Building a map from the walkthrough. \\
    Ask before departure only & The person may not be present once the rover is moving. & Asking when confidence drops mid-mission. \\
    Unclear confirmation means stay & A missed yes is cheap, a missed no is not. & Moving on a best guess. \\
    Vocabulary stated in the prompt and audited & Lets me measure how often the planner leaves it. & Silently correcting every plan. \\
    Vocabulary enforced on recovery plans & Nobody hears a recovery plan before it runs. & Trusting the prompt alone. \\
    Recovery only on failure & Keeps it distinct from clarification. & One mechanism for both. \\
    Unsure condition means ask & A wrong guessed branch looks like a correct one. & Defaulting to yes or no. \\
    Condition action offered only when needed & Keeps the prompt short and stops a condition being invented. & Listing every action in every prompt. \\
    A recovery plan must retry the approach & A plan that walks away from the failed step is not recovery. & Accepting any valid plan. \\
    Gimbal-first search & Pan angles are known, pivots drift. & Spinning the chassis to look. \\
    Two-signal arrival & Neither sonar nor a 3B vision model is reliable alone. & Trusting the first positive. \\
    Fixed hop as default, variable as option & The distance estimate is uncalibrated, and I want it tested first. & Committing to the variable hop. \\
    \bottomrule
  \end{tabular}
\end{table}
